\documentclass[11pt]{article}
\usepackage[margin=1in]{geometry}
\usepackage[T1]{fontenc}
\usepackage[utf8]{inputenc}
\usepackage{lmodern}
\usepackage{amsmath,amssymb,amsthm,mathtools}
\usepackage{booktabs,tabularx,longtable,array}
\usepackage{xltabular}% tabularx-style tables that can break across pages
\usepackage{enumitem}
\usepackage[dvipsnames]{xcolor}
\usepackage{cite}% sorts and compresses numeric citations
\usepackage[colorlinks=true,linkcolor=Blue,citecolor=Blue,urlcolor=Blue]{hyperref}

% ---------------------------------------------------------------------
% Shared notation
\newcommand{\R}{\mathbb{R}}
\newcommand{\E}{\mathbb{E}}
\renewcommand{\P}{\mathbb{P}}
\newcommand{\one}{\mathbf{1}}
\newcommand{\D}{\mathcal{D}}
\newcommand{\Lgen}{\mathcal{L}}
\DeclareMathOperator{\Var}{Var}
\DeclareMathOperator{\Cov}{Cov}
\DeclareMathOperator{\Ent}{Ent}
\DeclareMathOperator{\KL}{KL}
\DeclareMathOperator{\Tr}{Tr}
\DeclareMathOperator{\ran}{ran}
\newcommand{\norm}[1]{\lVert #1\rVert}
\newcommand{\ip}[2]{\langle #1,#2\rangle}
% Locators.  \ploc{compiled locator}{TeX file:line} cites the preprint (or a
% companion).  Only the compiled locator (section, equation, result, page of the
% PDF) is printed.  The TeX file:line stays in the source to help checkers find
% the passage and never appears in the PDF.  \sloc{key}{locator} cites a source.
\newcommand{\ploc}[2]{#1}
\newcommand{\sloc}[2]{\cite[#2]{#1}}
% Relationship types and flags, as used in the summary table and the Summary
% paragraphs: \rel{same technique}; \attrgap and \overstated tag a Cited?
% cell (precedence_stats.py checks them against the JSON).
\newcommand{\rel}[1]{\textit{#1}}
\newcommand{\attrgap}{\textbf{attribution gap}}
\newcommand{\overstated}{\textbf{overstated novelty}}
% Uniform list and table styling across all findings.
\setlist{itemsep=2pt,topsep=3pt,parsep=0pt,leftmargin=1.6em}
\setlist[enumerate]{label=\arabic*.}
\setlist[itemize]{label=\textbullet}
\renewcommand{\arraystretch}{1.15}
% Findings are numbered by LaTeX, so nobody types an ID.  In the main file,
% \findinggroup{F} after the Findings \section sets the ID prefix and restarts
% the count; IDs follow the order of the \input lines (by magnitude).  Each
% fragment opens with
%   \finding{<title>}{<slug>}   or   \finding[<short title>]{<title>}{<slug>}
% which prints an unnumbered subsection "F3. <title>", adds "F3. <title>" (or
% the short title) to the table of contents and PDF bookmarks, and does
% \label{f:<slug>}, so \ref{f:<slug>} prints the ID "F3" as a link.  The
% heading carries no LaTeX section number, so a finding never shows two
% numbers (no "3.2 F2").  Refer to another finding as finding~\ref{f:<slug>}.
% Before any \findinggroup the IDs print as ?1, ?2, ... (as in a standalone
% test of one fragment).
\makeatletter
\newcounter{finding}
\newcommand{\finding@prefix}{?}
\renewcommand{\thefinding}{\finding@prefix\arabic{finding}}
\newcommand{\findinggroup}[1]{\gdef\finding@prefix{#1}\setcounter{finding}{0}}
\newcommand{\finding}[3][]{%
  \stepcounter{finding}%
  \protected@edef\finding@id{\thefinding}% fix the ID now: "F3", not a counter reference
  \edef\finding@next{\noexpand\finding@sect{\finding@id}}%
  \finding@next{#1}{#2}{#3}}
% #1 ID (already expanded), #2 short title (may be empty), #3 title, #4 slug.
\newcommand{\finding@sect}[4]{%
  \subsection*{#1. #3}% hyperref anchors \subsection* itself, just above the heading
  \if\relax\detokenize{#2}\relax
    \addcontentsline{toc}{subsection}{#1. #3}%
  \else
    \addcontentsline{toc}{subsection}{#1. #2}%
  \fi
  \def\@currentlabel{#1}% \ref prints the ID; hyperref links it to this heading
  \label{f:#4}}
\makeatother
% ---------------------------------------------------------------------
\space%

\title{Missing citations in\\ ``Full support of the zero-temperature
Sherrington--Kirkpatrick order parameter''}
\date{October 13, 2026}

\begin{document}
\maketitle

\section*{Summary of findings}

Each row is a citation the preprint is missing: earlier work that already
contains part of the preprint, which the preprint does not cite where that
part appears. ``Cited?'' names the earlier work and says what, if anything, the
preprint cites it for. The exception is a row flagged \emph{overstated
novelty}: there the preprint cites the work, but its abstract or framing
presents the content as new. Terms are defined in Section~\ref{sec:about};
centrality is a best-effort judgment. Locations use the preprint's own
numbering and printed pages.

% Each row starts "\ref{f:<finding>} & \rel{<relationship>} & <centrality> &",
% in that order: \ref prints the finding's ID (F1, F2, ...) as a link, and
% precedence_stats.py reads those three cells. Put \attrgap and \overstated in
% the Cited? cell of a row exactly when the JSON sets the flag. Never type an
% ID: LaTeX assigns them from the main file's order (by magnitude).
% The Technique or result cell, like a title, names the approach or result,
% with no author names or citations.
\begin{center}\footnotesize
\begin{xltabular}{\linewidth}{@{}l >{\raggedright\arraybackslash}p{0.11\linewidth} >{\raggedright\arraybackslash}p{0.10\linewidth} >{\raggedright\arraybackslash}X >{\raggedright\arraybackslash}p{0.17\linewidth} >{\raggedright\arraybackslash}p{0.24\linewidth}@{}}
\toprule
ID & Relationship & Centrality & Technique or result & In the preprint & Cited?\\
\midrule
\endfirsthead
\toprule
ID & Relationship & Centrality & Technique or result & In the preprint & Cited?\\
\midrule
\endhead
\ref{f:chen-same-main-theorem} & \rel{same result} & main result &
Full support on $[0,1)$ of the zero-temperature order parameter, with strict
increase and (1.6) &
Abstract (p.~1); Thm~1.1 (p.~2) &
H.-B. Chen 2026 \cite{Chen2026-FRSBFullIntervalSupportZeroTemp}, first in
its aiXiv version \cite{ChatGPT2026-FRSBFullIntervalSupportZeroTemp}: cited
after Thm~1.1 (p.~2), not in the abstract; \overstated\\
\addlinespace
\ref{f:zhou-gap-exclusion} & \rel{same technique} & central &
Gap exclusion by the endpoint contact relations against a positive third
derivative on the gap &
\S1.2 and Fig.~1 (p.~4); Lemma~2.3 (p.~7); proof of Thm~1.1 (p.~20) &
Zhou 2025 \cite{Zhou2025-FRSBSKLowTemperature}: cited for other content;
Lopatto 2026 \cite{Lopatto2026-FullRSBSK}: Prop.~9.2 credited elsewhere
(\S1.1, p.~3), Lemma~2.5 cited for other content; H.-B. Chen 2026
\cite{Chen2026-FRSBFullIntervalSupportZeroTemp}: cited for other content;
\attrgap\\
\addlinespace
\ref{f:lopatto-slope-invariant} & \rel{same result} & supporting &
The Cole--Hopf slope invariant as sign inequalities for the backward drift &
Lemma~3.2 (p.~11) &
Lopatto 2026 \cite{Lopatto2026-FullRSBSK}: credited elsewhere (\S1.1,
p.~3); H.-B. Chen 2026 \cite{Chen2026-FRSBFullIntervalSupportZeroTemp}:
credited elsewhere in part\\
\addlinespace
\ref{f:zhou-ratio-sign} & \rel{same result} & supporting &
Monotonicity of the ratio of the forward score to the drift, invariant at
jumps &
Lemma~3.3 (pp.~12--13) &
Zhou 2026 \cite{Zhou2026-GardnerTransitionIsingPurePSpinII}: not cited;
\attrgap\\
\addlinespace
\ref{f:chen-terminal-boundary-layer} & \rel{same result} & supporting &
Exclusion of a terminal plateau by a square-root boundary layer &
Lemma~6.3 (pp.~19--20) &
H.-B. Chen 2026 \cite{Chen2026-FRSBFullIntervalSupportZeroTemp}, first in
its aiXiv version \cite{ChatGPT2026-FRSBFullIntervalSupportZeroTemp}: cited
for other content; Auffinger--Chen--Zeng 2017
\cite{AuffingerEtAl2017-SKInfiniteStepRSBZeroTemp}: credited for the
conclusion; El Alaoui--Montanari--Sellke 2020
\cite{ElAlaouiEtAl2020-OptimizationMeanFieldSpinGlasses}: credited in part;
\attrgap\\
\addlinespace
\ref{f:chl-origin-in-support} & \rel{same result} & supporting &
The origin lies in the support: no initial interval of zero mass &
Proof of Thm~1.1, initial-gap case (p.~20) &
Chen--Handschy--Lerman 2016 \cite{ChenEtAl2016-EnergyLandscapeMixedEvenPSpin}:
cited for other content; Auffinger--Chen 2013
\cite{AuffingerChen2013-PropertiesParisiMeasures}: cited for other content;
\attrgap\\
\addlinespace
\ref{f:chen-eams-hypothesis-verified} & \rel{same result} & peripheral &
Full support verifies the no-overlap-gap assumption, giving an
$\epsilon$-optimal algorithm &
Cor.~7.3 (p.~25); announced on pp.~2--3 &
H.-B. Chen 2026 \cite{Chen2026-FRSBFullIntervalSupportZeroTemp}: cited for
other content\\
\bottomrule
\end{xltabular}
\end{center}

\clearpage
\space%

\tableofcontents

\section{About this document}\label{sec:about}

\paragraph{The preprint.}
The preprint is \emph{Full support of the zero-temperature
Sherrington--Kirkpatrick order parameter}, dated September~27, 2026, with
author line ``OpenAI'' \cite{OAI-Preprint}. It is one of the preprints of the
OpenAI Math Release. It proves that every minimizer of the zero-temperature
Parisi functional for the pure, zero-field Sherrington--Kirkpatrick model has
full support on $[0,1)$, and that the minimizing order parameter has a smooth,
strictly positive density on $(0,1)$. Section~\ref{sec:overview} summarizes
the paper and fixes the notation used throughout this document. The preprint
cites no other preprint from the same release.

\paragraph{What this document reports.}
Each finding is a citation the preprint is missing: part of the preprint's
mathematics (a technique, construction, estimate, proof route or result) that
an identifiable earlier paper already contains, where the preprint does not
cite that paper for it where it is used. ``Earlier'' means publicly available
before the preprint's date. The one exception, overstated novelty, is
described below. Not reported: textbook material in its textbook form; weaker versions of
intermediate statements; and, for content the preprint uses, cases where it
cites a later source that covers that content rather than the original. An
earlier result on the same problem is reported even when later work that the
preprint cites implies it. Findings are ordered by centrality, then by relationship in
the order listed below, then by position in the preprint.

\paragraph{Relationships.}
A finding's relationship is the one between the preprint and the first
earlier work listed for it (the principal), the work with the highest relationship in the
order below; when several hold that relationship, however closely each matches,
the one the others credit or build on, otherwise the earliest. Any other works listed relate to the same content more weakly or
in part; the finding says how, including which work originated the idea when
that is not the first.
\begin{description}[leftmargin=1.5em,font=\normalfont\itshape]
\item[same result.] An earlier paper, or several earlier papers together, prove essentially the same statement.
\item[same technique.] The earlier technique, construction or theory is
reused essentially intact, possibly in a new setting or with new ingredients
around it.
\item[adapted.] The earlier technique's structure is followed, with a change
at its core: a different object, or a key step replaced.
\item[proof shortcut.] A published result would replace or substantially
shorten a large part of the argument.
\item[same route.] An earlier paper proposed the proof strategy, even
heuristically, and the preprint carries it out.
\item[fragment.] One identifiable piece of the earlier technique is reused.
\item[weaker result.] An earlier paper proves a weaker, special-case or
partial-range version of one of the preprint's main results, or of a
statement it presents as a significant contribution.
\end{description}

\paragraph{Centrality.}
Each finding also gets a best-effort judgment of how much of the preprint's
main argument depends on it: \emph{main result} (reserved for earlier work that
covers the main result itself, or most of it), \emph{central},
\emph{supporting} or \emph{peripheral}. An earlier weaker result is never
labelled main result. Each finding states the chain of results it feeds, so
the basis for the judgment is visible.

\paragraph{Citations.}
The ``Cited?'' column of the summary names the earlier work and says what the
preprint cites for the content in question: \emph{not cited} (the work is not
in the bibliography, and nothing the preprint cites credits it for this
content), \emph{cited for other content}, \emph{credited elsewhere} (credited
for this content, directly or through a cited location that credits it, but
not where it is used, in a way that an expert reader who doesn't know the earlier work might not connect with it), or \emph{credited in part}. A citation counts only when it is \emph{for that specific content};
citing the same paper nearby for a different result does not count. Content
the preprint credits where it is used, directly or through a cited location
that itself credits the earlier work, is not reported, except when its novelty is overstated (below). For
overstated novelty, the column says where the preprint does cite the work.

\paragraph{Flags.}
\begin{description}[leftmargin=1.5em,font=\normalfont\bfseries]
\item[Attribution gap.] Marks the most serious missing citations: the
relationship is same technique, adapted or same result; the centrality is main
result, central or supporting; and the preprint cites the principal earlier
work for this content nowhere.
\item[Overstated novelty.] The preprint does cite the earlier work for this
content, but its abstract, a statement of its significant contributions, or the framing of
a main result presents that content as new, explicitly or in the way a reader
would reasonably take it. This is the one case in which content the preprint
credits is reported.
\end{description}

\paragraph{How to check a claim.}
Locations in the preprint use its section,
equation and theorem numbers and the page numbers printed in its PDF.
Locations in other sources give the theorem, equation and printed page of the
version listed in the bibliography; where an entry lists both a journal and an
arXiv version, they refer to the arXiv version. Each finding ends with a short checklist
of the passages to put side by side.

\paragraph{Provenance.}
This document was generated by LLM agents (Claude Opus 5.5) running the
precedence-review skill
(\href{https://github.com/jnshi/math-lineage/tree/13ccc19ec/.agents/skills/precedence-review}{jnshi/math-lineage},
commit \texttt{13ccc19}; parts of the review ran under the earlier versions
\texttt{d1ea740} and \texttt{3808a6c}), makes no claim of completeness, and
has not been vetted by a human. It was run by Jonathan Shi.
\space%
\section{The preprint in brief, and notation}\label{sec:overview}

\paragraph{Model and variational problem.}
The pure, zero-field Sherrington--Kirkpatrick Hamiltonian is
$H_n(\sigma)=n^{-1/2}\sum_{i<j}J_{ij}\sigma_i\sigma_j$, normalized so that the
covariance is $\xi(t)=t^2/2$ and $\xi''\equiv1$. The class $\mathcal U$
consists of nonnegative, nondecreasing, right-continuous
$\gamma:[0,1)\to[0,\infty)$ with $\int_0^1\gamma<\infty$, and $\mu_\gamma$ is
the Stieltjes measure with $\mu_\gamma([0,t])=\gamma(t)$. The Parisi value
function $\Phi_\gamma$ solves
\[
 \partial_t\Phi=-\tfrac12\bigl(\Phi_{xx}+\gamma(t)\Phi_x^2\bigr),
 \qquad \Phi(1,x)=|x|,
\]
\ploc{(1.1)--(1.2)}{main.tex:65--78}, and the zero-temperature Parisi
functional is $\mathcal P(\gamma)=\Phi_\gamma(0,0)-\frac12\int_0^1t\gamma(t)\,dt$
\ploc{(1.3)}{main.tex:82--85}. Its minimum over $\mathcal U$ is the limiting
ground-state energy per spin \ploc{(1.4)}{main.tex:100--103}.

\paragraph{The optimal diffusion and the self-consistency function.}
Write $u=\Phi_{\gamma,x}$ and let $X$ solve
$dX_t=\gamma(t)u(t,X_t)\,dt+dB_t$, $X_0=0$ \ploc{(1.5)}{main.tex:109--113}. Set
\[
 q(t)=\E\bigl[u(t,X_t)^2\bigr],\qquad q'(t)=\E\bigl[\Phi_{\gamma,xx}(t,X_t)^2\bigr]
\]
\ploc{(2.12), Lemma~2.2}{regularity.tex:189--208}. The sources below write
$\Gamma$ for this function, and $m$ for the value of a constant stretch of the
order parameter.

\paragraph{Main theorem.}
\ploc{Theorem~1.1}{main.tex:119--131}: if $\gamma\in\mathcal U$ minimizes
$\mathcal P$, then $\operatorname{supp}\mu_\gamma=[0,1)$, and $q(t)=t$ and
$\E[\Phi_{\gamma,xx}(t,X_t)^2]=1$ for $0\le t<1$. Corollary~6.4 adds that
$\gamma$ is smooth on $[0,1)$ with $\gamma(0)=0$ and $\gamma'>0$ on $(0,1)$.

\paragraph{Contact conditions.}
For a minimizer, the contact potential
$G(s)=\int_s^1(q(t)-t)\,dt$ is nonnegative and vanishes on the support; at
every interior support point, $q(s)=s$ and $q'(s)\le1$
(\ploc{Lemma~6.2, p.~19}{variation.tex:81--128}).

\paragraph{Notation on a constant interval.}
On an interval where $\gamma\equiv c>0$, the preprint sets
\[
 r=c\Phi_x,\qquad v=r_x,\qquad w=r_{xx},\qquad z=r_{xxx},
\]
and proves
\[
 (q')'=c^{-2}\E\bigl[w^2-2v^3\bigr],\qquad
 (q')''=c^{-2}\E\bigl[z^2-12vw^2+6v^4\bigr]
\]
along the diffusion (\ploc{(2.13)--(2.14), Lemma~2.3}{regularity.tex:240--261}).

\paragraph{Proof route.}
\begin{itemize}
\item \ploc{Proposition~5.1}{curvature.tex:9--20}: on every positive constant
interval, $q'''\ge(1000c^2)^{-1}\E v^4>0$, so $q'$ is strictly convex there.
\item Interior gaps (\ploc{proof of Theorem~1.1, p.~20}{variation.tex:215--250}):
on a gap $(a,b)$ with endpoints in the support, $q(a)=a$, $q(b)=b$ and
$q'(a),q'(b)\le1$. Strict convexity puts $q'$ below its chord, hence below $1$,
on $(a,b)$, so $q(b)-q(a)<b-a$, a contradiction
(\ploc{Figure~1, p.~4}{gap-convexity.tex}). Gaps on which $\gamma=0$ are
excluded by strict monotonicity of $q'$.
\item A terminal plateau is excluded separately
(\ploc{Lemma~6.3, pp.~19--20}{variation.tex:130--213}).
\end{itemize}

\paragraph{The machinery behind Proposition~5.1.}
The preprint's \S\S3--4 work with a positive finite-step coefficient
$c_1\le\dots\le c_N\le M$ and terminal datum $M^{-1}\log(2\cosh Mx)$
(\ploc{\S3, p.~8}{signs.tex:17--23}). They factor the density of $X_t$ as
$f=e^{c\Phi}g$, where $g$ solves the forward heat equation between jumps, and
put $s=-\partial_x\log g$ (\ploc{\S3, p.~8}{signs.tex:35--56}). With backward
time $\tau$ and $H=r_\tau/r$:
\begin{itemize}
\item \ploc{Lemma~3.2}{signs.tex:195--203}: on $x>0$, $w\le0$, $H\le0$ and
$H_x\ge0$;
\item \ploc{Lemma~3.3}{signs.tex:289--297}: $s_x\ge1/t$, $s_{xx}\le0$ on $x>0$,
and $j:=s_x-(s/r)v=r\,\partial_x(s/r)\ge0$;
\item \ploc{\S3.4 and Lemma~3.4}{signs.tex:396--433}: the variables
$R=r^2$, $K=rs$, $b=-w/r$, $B=b-2v=-2H$, $S=z+bv-2Rb=2rH_x$ and
$L=-rs_{xx}$ are nonnegative, and
$\E_f[\D F+(v+R-K)F]=0$ for polynomials $F$, where $\D=r\partial_x$;
\item \ploc{Lemma~4.1 and Proposition~4.2}{certificate.tex:117--193}: an exact
identity over $\mathbb Q$ writes $z^2-12vw^2+6v^4$ as a sum of nonnegative
terms plus such a divergence, giving
$\E_f[z^2-12vw^2+6v^4]\ge10^{-3}\E_fv^4$.
\end{itemize}
\space%

% IDs are assigned by LaTeX in the order of the \input lines (F1, F2, ...).
% Order by magnitude: centrality (main result, central, supporting,
% peripheral), then the relationship (same result, same technique, adapted,
% proof shortcut, same route, fragment, weaker result), then reading order.
% precedence_stats.py checks this order.
% Withdrawing or reordering a finding is just deleting or moving its \input line.
\section{Findings}\label{sec:findings}
\findinggroup{F}
\finding{Full support of the zero-temperature order parameter}{chen-same-main-theorem}

\paragraph{Summary.}
H.-B.~Chen's paper \sloc{Chen2026-FRSBFullIntervalSupportZeroTemp}{Thm~1.3,
p.~3} (arXiv, July~20, 2026) proves that the minimizer of the
zero-temperature Parisi functional of the pure, zero-field SK model has full
support on $[0,1)$; the theorem first appeared in its aiXiv version, posted
July~16, 2026 \sloc{ChatGPT2026-FRSBFullIntervalSupportZeroTemp}{Thm~1.2,
p.~2}, which Chen's paper credits
\sloc{Chen2026-FRSBFullIntervalSupportZeroTemp}{p.~1; App.~C, p.~38}.
The preprint cites Chen's paper for this theorem right after
\ploc{Theorem~1.1, p.~2}{main.tex:133--135}, but its abstract
(\ploc{p.~1}{main.tex:33--36}) opens ``We prove that'' and presents the
theorem as new, naming neither Chen's paper nor an earlier proof.
The relationship is \rel{same result}: every assertion of
\ploc{Theorem~1.1}{main.tex:119--131} (full support, strict increase of
$\gamma$, and \ploc{(1.6)}{main.tex:125--130}) is in Chen's paper, for the
same model, class, functional and notion of support.
Both proofs follow the plan of Chen's paper, first set out in its aiXiv
version \sloc{ChatGPT2026-FRSBFullIntervalSupportZeroTemp}{pp.~2--3}: exclude
each component of the complement of the support directly at zero
temperature, using that $\gamma$ is constant on it and the first-variation
contact conditions (a gap between support points through the frame of
finding~\ref{f:zhou-gap-exclusion}, a terminal plateau through the boundary
layer of finding~\ref{f:chen-terminal-boundary-layer}).
In the gap step, Chen's paper gets the contradiction from the crossing rule
$\Gamma''=0\Rightarrow\Gamma'''>0$
\sloc{Chen2026-FRSBFullIntervalSupportZeroTemp}{Prop.~4.2, p.~19} and the
preprint from the unconditional bound $q'''\ge(1000c^2)^{-1}\E v^4$ of
\ploc{Proposition~5.1}{curvature.tex:9--20}, a difference the preprint states
(\ploc{pp.~2--3}{main.tex:135--137, 198--203}).

\paragraph{Correspondence.}
\begin{center}\small
\begin{xltabular}{\linewidth}{@{}>{\raggedright\arraybackslash}p{0.15\linewidth}>{\raggedright\arraybackslash}X>{\raggedright\arraybackslash}X@{}}
\toprule
 & Preprint & Chen \cite{Chen2026-FRSBFullIntervalSupportZeroTemp}\\
\midrule
\endfirsthead
\toprule
 & Preprint & Chen \cite{Chen2026-FRSBFullIntervalSupportZeroTemp}\\
\midrule
\endhead
Setting & $\xi(t)=t^2/2$; $\mathcal U$, $\mu_\gamma([0,t])=\gamma(t)$; (1.2)--(1.3) (\ploc{pp.~1--2}{main.tex:41--87}) & $\xi(t)=t^2/2$; $\mathcal U$ (1.4), $\mathrm d\gamma([0,t])=\gamma(t)$; (1.5)--(1.6) (p.~3)\\
\addlinespace
Full support & every minimizer $\gamma\in\mathcal U$: $\operatorname{supp}\mu_\gamma=[0,1)$, and $\gamma(b)>\gamma(a)$ for $0\le a<b<1$ (\ploc{Theorem~1.1, p.~2}{main.tex:120--124}) & the minimizer $\gamma_\star$, unique by Chen, Handschy and Lerman \sloc{ChenEtAl2016-EnergyLandscapeMixedEvenPSpin}{Thm~4, p.~8}: $\operatorname{supp}_{[0,1)}\nu_\star=[0,1)$, $\nu_\star=\mathrm d\gamma_\star$ (Thm~1.3, (1.8), p.~3), and $\gamma_\star(t)-\gamma_\star(s)=\nu_\star((s,t])>0$ for $0\le s<t<1$ (p.~4)\\
\addlinespace
Saturated consistency & $\E[u_\gamma(t,X_t)^2]=t$, $\E[\Phi_{\gamma,xx}(t,X_t)^2]=1$, $0\le t<1$ \ploc{(1.6)}{main.tex:125--130} & $\Gamma(t)=t$ (4.40), $\E\,\mathsf C(t,X_t)^2=1$ (4.41), $0\le t<1$, where $\mathsf C=u_{xx}$ and $u=\Phi_\gamma$ (proof of Prop.~4.6, p.~25)\\
\addlinespace
\bottomrule
\end{xltabular}
\end{center}

\paragraph{How to check.}
\begin{enumerate}
\item \ploc{Theorem~1.1, p.~2}{main.tex:119--131} against
\sloc{Chen2026-FRSBFullIntervalSupportZeroTemp}{Thm~1.3, (1.8), p.~3}, the
display on p.~4 of Chen's paper, and
\sloc{ChatGPT2026-FRSBFullIntervalSupportZeroTemp}{Thm~1.2, p.~2}, whose
(1.3) is the second identity in Chen's (1.8).
\item \ploc{(1.6)}{main.tex:125--130} against
\sloc{Chen2026-FRSBFullIntervalSupportZeroTemp}{(4.40)--(4.41), p.~25}.
\item The preprint's abstract (\ploc{p.~1}{main.tex:33--36}) against its
citation of Chen's paper after \ploc{Theorem~1.1, p.~2}{main.tex:133--135}.
\end{enumerate}

\space%
\finding[Gap exclusion by endpoint contact and a positive third derivative]{Excluding support gaps by the endpoint contact relations and a positive third derivative}{zhou-gap-exclusion}

\paragraph{Summary.}
Zhou excludes a gap $(a,b)$ near the origin in the support of the
near-critical SK Parisi measure by showing that the endpoint relations
$\Gamma(a)=a$, $\Gamma(b)=b$, $\Gamma'(a),\Gamma'(b)\le1$
\sloc{Zhou2025-FRSBSKLowTemperature}{(7), p.~9} contradict $G(a,b)>0$, which
for short gaps he derives from a lower bound on
$\partial_t^2G(a,a+)=\Gamma'''(a+)/(3\Gamma'(a))$ and a bound on
$\partial_t^3G$ \sloc{Zhou2025-FRSBSKLowTemperature}{Thm~5(2)--(3), Cor.~6,
p.~8; pp.~14--15}. The preprint cites Zhou
only for his interval-support theorem
(\ploc{\S1.1, p.~3}{main.tex:189--190}), and not where it uses his
approach: \ploc{\S1.2, p.~4}{main.tex:243--257; gap-convexity.tex:1--31},
\ploc{Lemma~2.3, p.~7}{regularity.tex:240--261} (whose (2.14) is Zhou's
$\Gamma'''$ formula and sets the target of
\ploc{\S\S3--4}{signs.tex:4--6; certificate.tex:3--7}) and the
\ploc{proof of Theorem~1.1, p.~20}{variation.tex:223--234} (its initial-gap
case is finding~\ref{f:chl-origin-in-support}).
The approach and its insight are Zhou's (Lopatto
\sloc{Lopatto2026-FullRSBSK}{p.~3} calls the general strategy closely related
to Zhou's argument): on a gap the order parameter is constant, so $\Gamma$ has
explicit higher derivatives along the optimal diffusion, and
$\Gamma'''>0$ makes the trapezoid value $(b-a)(\Gamma'(a)+\Gamma'(b))/2$,
at most $b-a$ by the slope relations, exceed $\int_a^b\Gamma'=b-a$.
\rel{same technique}: the frame is reused intact, and the integrated chord
bound is Zhou's $G(a,b)>0$; the steps differ only in the positivity the frame
needs, which \ploc{Proposition~5.1, p.~17}{curvature.tex:9--20} gives on every
positive constant interval, where Zhou's bound is local and near-critical
and the crossings of Lopatto and Chen 2026, which the preprint
cites, are conditional.
Lopatto has the same formulas \sloc{Lopatto2026-FullRSBSK}{Lemma~2.5, p.~6},
never cited, and the frame with such a crossing
\sloc{Lopatto2026-FullRSBSK}{Prop.~9.2, pp.~50--51}, cited only through his
Sections~3, 7--9 in the introduction (\ploc{p.~3}{main.tex:192--195});
Chen's paper \sloc{Chen2026-FRSBFullIntervalSupportZeroTemp}{Prop.~4.3,
pp.~22--23}, cited for other results, adapts the
frame at zero temperature, closing with a weighted identity instead of the
chord (first in its aiXiv version of July~16, 2026
\sloc{ChatGPT2026-FRSBFullIntervalSupportZeroTemp}{Prop.~4.2, pp.~30--31}).

\paragraph{Correspondence.}
\begin{center}\small
\begin{xltabular}{\linewidth}{@{}>{\raggedright\arraybackslash}p{0.11\linewidth}>{\raggedright\arraybackslash}p{0.28\linewidth}>{\raggedright\arraybackslash}X@{}}
\toprule
 & Preprint & Earlier work\\
\midrule
\endfirsthead
\toprule
 & Preprint & Earlier work\\
\midrule
\endhead
Endpoint relations & $q(a)=a$, $q(b)=b$, $q'(a),q'(b)\le1$ (\ploc{(6.4), p.~19}{variation.tex:94--99}; \ploc{p.~20}{variation.tex:223--226}) & Zhou (7), p.~9. Lopatto (9.32), p.~51: $h(a)=h(b)=0$, $h'(a),h'(b)\le0$, $h=\Gamma-s$. Chen 2026 (4.23)--(4.24), p.~22: $\Gamma(a)=a$, $\Gamma(b)=b$, $\Gamma'(a+)\le1$, $\int_a^b(\Gamma-t)=0$\\
\addlinespace
Third derivative & $(q')''=c^{-2}\E[z^2-12vw^2+6v^4]$ (\ploc{(2.13)--(2.14), p.~7}{regularity.tex:254--258}) & Zhou p.~15, with $8\beta^6=(\xi'')^3$, here $1$, and $m=c$: $\Gamma'''=8\beta^6\E[((\partial_x^4\Phi)^2-12m\,\partial_x^2\Phi(\partial_x^3\Phi)^2+6m^2(\partial_x^2\Phi)^4)e^{W}]$. Lopatto (2.18), p.~6. Chen 2026 (3.38), p.~14: $\Gamma''$ only\\
\addlinespace
Closing step & $q'$ below its chord, hence below $1$, on $(a,b)$; integrating, $q(b)-q(a)<b-a$ (\ploc{p.~20}{variation.tex:227--234}) & Zhou: $G(a,b)>0$, $G(s,t)=(t-s)[\Gamma'(s)+\Gamma'(t)]/(\Gamma(t)-\Gamma(s))-2$ (pp.~8--9). Lopatto: $\max h'=\max\{h'(a),h'(b)\}\le0$, $0=\int_a^bh'$ (p.~51). Chen 2026: $\int_a^b(t-a)(b-t)\Gamma''=0$, so $\Gamma(b)-\Gamma(a)<(b-a)\Gamma'(a+)$ ((4.25)--(4.26), p.~23)\\
\addlinespace
\bottomrule
\end{xltabular}
\end{center}

\paragraph{How to check.}
\begin{enumerate}
\item Zhou's $G$ (p.~8) and (7) (p.~9) against the chord display in the
\ploc{proof of Theorem~1.1, p.~20}{variation.tex:229--234}: integrated, it
gives $G(a,b)>0$ for $\Gamma=q$, while the four relations give $G(a,b)\le0$.
\item Zhou's $\Gamma'''$ (p.~15) and Lopatto's Lemma~2.5 (p.~6) against
\ploc{Lemma~2.3 and its proof, pp.~7--8}{regularity.tex:240--296}.
\item Chen 2026, Prop.~4.3 (pp.~22--23; aiXiv, Prop.~4.2) against the
\ploc{proof of Theorem~1.1, p.~20}{variation.tex:223--234}.
\end{enumerate}

\space%
\finding[Cole--Hopf slope invariant as backward drift signs]{The Cole--Hopf slope invariant as sign inequalities for the backward drift}{lopatto-slope-invariant}

\paragraph{Summary.}
Lopatto's Proposition~3.6, which he calls the Cole--Hopf slope invariant,
shows that every finite Cole--Hopf cascade from $\log\cosh$ satisfies
$K\ge0$ and four further inequalities, where $KB=x_r$ is the speed of the
level sets of the slope $B=u_x$ \sloc{Lopatto2026-FullRSBSK}{pp.~15--18};
Chen's paper transports them to zero temperature
by smoothing $|x|$ to $\lambda^{-1}\log\cosh\lambda x$ and rescaling (proof of
Prop.~3.11), and proves $u_{xxx}\le0$ by a maximum principle (Lemma~3.4)
\cite{Chen2026-FRSBFullIntervalSupportZeroTemp}, both already in its aiXiv
version \cite{ChatGPT2026-FRSBFullIntervalSupportZeroTemp}. The preprint cites
neither where it uses this content: it mentions Lopatto's invariant only in
its introduction, in describing his proof, and cites Chen's Proposition~3.11
only there, for ``smooth-terminal approximation''
(\ploc{\S1.1, p.~3}{main.tex:192--197}); its \S3 cites Lopatto only for the
density transform (\ploc{p.~12}{signs.tex:285--286}), and Chen's Lemma~3.4
nowhere.
\ploc{Lemma~3.2, p.~11}{signs.tex:195--203} states $w\le0$, $H=r_\tau/r\le0$
and $H_x\ge0$ on $x>0$: \rel{same result},
because in Lopatto's variables (the table below; not the preprint's
certificate $B$, $K$) $H=-CK$, $H_x=-C(CK)_B$ and
$w=cCC_B=-2cC(c+K)B$, so $H\le0$ is his $K\ge0$, $w\le0$ follows from it (it
is his $C_B\le0$), and $H_x\ge0$ follows from his inequalities in a few lines.
The approach is Lopatto's: follow the sign of $K$ through a finite cascade
started at $\log\cosh$, where $K=0$, using that each Cole--Hopf step
preserves it (by a maximum principle for $J=(BK)_B$ and the average (3.61))
and that lowering the parameter from $a$ to $b$ adds
$a-b$ to $K$; the preprint's maximum principles for $H$ and $H_x$ in $x$ are
its own, and its proof of $w\le0$ follows Chen's paper.

\paragraph{Correspondence.}
\begin{center}\small
\begin{xltabular}{\linewidth}{@{}>{\raggedright\arraybackslash}p{0.07\linewidth}>{\raggedright\arraybackslash}X>{\raggedright\arraybackslash}X@{}}
\toprule
 & Preprint, \S3 & Lopatto \cite{Lopatto2026-FullRSBSK}; Chen \cite{Chen2026-FRSBFullIntervalSupportZeroTemp}\\
\midrule
\endfirsthead
\toprule
 & Preprint, \S3 & Lopatto \cite{Lopatto2026-FullRSBSK}; Chen \cite{Chen2026-FRSBFullIntervalSupportZeroTemp}\\
\midrule
\endhead
Data & $0<c_1\le\dots\le c_N\le M$, terminal $M^{-1}\log(2\cosh Mx)$
(\ploc{p.~8}{signs.tex:17--23}); on a step, $\gamma=c$ & cascades from
$(\log\cosh,1)$ (Prop.~3.6, p.~15); $a=c$, $u=\Phi$, $B=r/c$, $C=v/c$, his
time $r$ is $\tau$; Chen (3.47), p.~15: after $y=\lambda x$, coefficients
$\le\lambda$ from $\lambda^{-1}\log\cosh\lambda x$ give such cascades\\
\addlinespace
Drift & $H=r_\tau/r=w/(2r)+v$ \ploc{(3.8), p.~11}{signs.tex:197--200} &
$x_r=KB$ (3.44), p.~16, and $C_B=-2(a+K)B$ (3.39), p.~15, so $H=-CK$\\
\addlinespace
Signs & $w\le0$, $H\le0$, $H_x\ge0$ on $x>0$ \ploc{(3.8)}{signs.tex:197--200};
$w\le0$ by a maximum principle (\ploc{p.~11}{signs.tex:220--229}) & $K\ge0$ (3.41), p.~16; $C_B\le0$ (p.~17);
Chen: $u_{xxx}\le0$ (3.20), p.~10, by the same maximum principle, (B.11),
p.~32\\
\addlinespace
Steps & terminal $H=(cM-M^2)\operatorname{sech}^2(Mx)$; jumps
$H_{\rm new}=H_{\rm old}-(1-\alpha)v_{\rm old}$,
$(H_x)_{\rm new}=(H_x)_{\rm old}-(1-\alpha)w_{\rm old}$, $\alpha=c_-/c_+$
(\ploc{pp.~11--12}{signs.tex:241--245, 270--275}) & $K=J=0$ at the start (3.67);
$K_b=K_a+a-b$ (3.65), p.~18; lowering $M$ to $c$ gives $K=M-c$, and (3.65)
gives the jump rules\\
\addlinespace
\bottomrule
\end{xltabular}
\end{center}

\paragraph{How to check.}
\begin{enumerate}
\item Put $r=cB$, $v=cC$, $w=cCC_B$ into \ploc{(3.8)}{signs.tex:197--200}
and use \sloc{Lopatto2026-FullRSBSK}{(3.38), (3.39), (3.44), pp.~15--16} to
get $H=-CK$. For $H_x\ge0$: with his $z=-\frac12C_B$ (3.42),
$\varphi:=-B(CK)_B$ vanishes at $B=0$, and
$\varphi_B={BK(a+2J-K)}+{zBK_B}+{CK_B}+{(3zJ-CJ_B)}\ge0$ by (3.41), (3.43).
\item Terminal $H$ and jumps (\ploc{pp.~11--12}{signs.tex:241--245, 270--275})
against \sloc{Lopatto2026-FullRSBSK}{(3.65), (3.67), p.~18}, scaled by
\sloc{Chen2026-FRSBFullIntervalSupportZeroTemp}{(3.47), p.~15}.
\item The proof of $w\le0$ on \ploc{p.~11}{signs.tex:220--229} against
\sloc{Chen2026-FRSBFullIntervalSupportZeroTemp}{Lemma~3.4, p.~10; App.~B.3,
(B.11), p.~32}, times $c$.
\end{enumerate}

\space%
\finding[Monotonicity of the ratio of score to drift]{Monotonicity of the ratio of the forward score to the drift}{zhou-ratio-sign}

\paragraph{Summary.}
Zhou's August 2026 paper on the Gardner transition in the Ising pure $p$-spin
glass proves in App.~E.3 that $w_0=W_y/B$ (in the table's variable $y$,
Zhou's $x$), the lower-history part of his density score divided by the slope
$B$, is nondecreasing on $y>0$, since $R=W_{yy}/C$ is
\sloc{Zhou2026-GardnerTransitionIsingPurePSpinII}{(E.37), (E.41), p.~128};
the preprint does not cite this paper, and cites Zhou only for interval
support near the critical temperature \cite{Zhou2025-FRSBSKLowTemperature}
(\ploc{\S1.1, p.~3}{main.tex:189--190}).
It states the ratio sign $j=r\,\partial_x(s/r)\ge0$, with its invariance at
forward jumps, in \ploc{Lemma~3.3, p.~12}{signs.tex:289--297}, as ``the
additional ratio sign needed by the coercive identity''
(\ploc{\S3.3}{signs.tex:287}), and proves it on
\ploc{p.~13}{signs.tex:338--393} from $a\ge1/t$ and $s_{xx}\le0$ (Zhou's
$W_{yy}\ge0$ and $-W_{yyy}\ge0$, (E.31)) and the backward signs of
\ploc{Lemma~3.2}{signs.tex:195--203}.
\rel{same result}: after the table's rescaling, $s/r=y/(\theta mB)+w_0/m$, and
$B>yC$ \sloc{Zhou2026-GardnerTransitionIsingPurePSpinII}{p.~130} makes the
first term nondecreasing.
Both need such a sign for forward-density terms of a third-derivative identity
on a constant interval (for Zhou, the covariance identity (E.27), which the
sign enters through (E.45); in the preprint,
Lemma~3.3~$\to$ \ploc{Lemma~3.4}{signs.tex:416--433}~$\to$
\ploc{Lemma~4.1 and Proposition~4.2}{certificate.tex:117--193}~$\to$
\ploc{Proposition~5.1}{curvature.tex:9--20}~$\to$
\ploc{Theorem~1.1}{main.tex:119--131}), and both use Zhou's insight of
measuring the score against the drift: a jump adds a multiple of the drift to
the score, so the ratio's monotonicity survives jumps, and on the Gaussian
first step it is concavity of the drift.
Between jumps Zhou propagates $R_y\ge0$ and averages, while the preprint
propagates $\rho=s_x/s-v/r\ge0$, each by a maximum principle; Zhou's p.~3 says
his appendix proofs were drafted by language models, pending his revision.

\paragraph{Correspondence.}
\begin{center}\small
\begin{xltabular}{\linewidth}{@{}>{\raggedright\arraybackslash}p{0.1\linewidth}>{\raggedright\arraybackslash\hsize=0.8\hsize}X>{\raggedright\arraybackslash\hsize=1.2\hsize}X@{}}
\toprule
 & Preprint \cite{OAI-Preprint} & Zhou \sloc{Zhou2026-GardnerTransitionIsingPurePSpinII}{App.~E, pp.~124--130}\\
\midrule
\endfirsthead
\toprule
 & Preprint \cite{OAI-Preprint} & Zhou \sloc{Zhou2026-GardnerTransitionIsingPurePSpinII}{App.~E, pp.~124--130}\\
\midrule
\endhead
Data & $0<c_1\le\dots\le c_N\le M$; terminal datum $M^{-1}\log(2\cosh Mx)$;
$X_0=0$ (\ploc{\S3, pp.~8--9}{signs.tex:17--56}) & masses
$m=c/M\in(0,1]$, atoms $d=\delta/M$; $W=0$ before the first atom (E.28);
upper cone (E.52) from $\log\cosh y$ (E.46), $J=1-m\ge0$ (E.50), also
for $m=1$\\
\addlinespace
Dictionary & $r=c\Phi_x$, $v=r_x$; $f=e^{c\Phi}g$, with $g$ a heat solution
between jumps and $g_+=g_-e^{-\delta\Phi}$ at a jump $\delta=c_+-c_-$
\ploc{(3.2), p.~9}{signs.tex:35--55}; $s=-\partial_x\log g$, $a=s_x$ &
$\theta=M^2t$, $y=Mx$ (Zhou's $t$, $x$); $U=M\Phi-\log2$ solves (E.1); $B=U_y$, $C=U_{yy}$
(E.2), not the $B$ of Section~\ref{sec:overview};
$g\propto pe^{-mU}\propto e^{-y^2/(2\theta)-W}$ (E.4)--(E.5), $p$ the density
of $MX_t$; $r=MmB$, $v=M^2mC$,
$s=x/t+MW_y$, $a=1/t+M^2W_{yy}$\\
\addlinespace
Ratio sign & $j=a-(s/r)v=r\,\partial_x(s/r)\ge0$
\ploc{(3.11), p.~12}{signs.tex:291--294}; first step: $s=x/t$ and $r\ge xv$;
later steps: maximum principle for $\rho=s_x/s-v/r$
\ploc{(3.16)}{signs.tex:338--376} & $s/r=y/(\theta mB)+w_0/m$, so
$j=M^2[(B-yC)/(\theta B)+C(R-w_0)]$ on $y>0$, with $B>yC$ (p.~130) and
$R\ge w_0$ (E.41), from $R_y\ge0$ (E.37), which uses (E.52); first step:
$W=0$\\
\addlinespace
Jumps & $j_+=j_-$ by cancellation (\ploc{p.~13}{signs.tex:378--390}), with
$s_+=s_-+\delta\Phi_x$ \ploc{(3.15), p.~12}{signs.tex:313--317} & the atom $W\mapsto W+dU$ (E.28) gives
$R\mapsto R+d$ and $w_0\mapsto w_0+d$; $B$, $C$ continuous\\
\bottomrule
\end{xltabular}
\end{center}

\paragraph{How to check.}
\begin{enumerate}
\item Check the dictionary (\ploc{\S3, pp.~8--9}{signs.tex:17--56}; Zhou's
  (E.1)--(E.5), (E.28)
  \sloc{Zhou2026-GardnerTransitionIsingPurePSpinII}{pp.~124--127}), derive
  $j$ from \ploc{(3.11)}{signs.tex:291--294}, (E.32) and (E.40), and read off
  its signs from (E.41) and $B>yC$ (p.~130).
\item Compare the first step and jump on
  \ploc{p.~13}{signs.tex:344--346, 378--390} with (E.28) and the line after
  (E.36) \sloc{Zhou2026-GardnerTransitionIsingPurePSpinII}{pp.~127--128}.
\end{enumerate}

\space%
\finding{Excluding a terminal plateau by a square-root boundary layer}{chen-terminal-boundary-layer}

\paragraph{Summary.}
H.-B.~Chen's paper \sloc{Chen2026-FRSBFullIntervalSupportZeroTemp}{Prop.~4.5,
pp.~23--25} proves that the zero-temperature minimizer has no constant stretch
ending at~$1$, first in its aiXiv predecessor of July~16, 2026
\sloc{ChatGPT2026-FRSBFullIntervalSupportZeroTemp}{Prop.~4.4, pp.~32--33},
which it credits
\sloc{Chen2026-FRSBFullIntervalSupportZeroTemp}{p.~1; App.~C, p.~38}; the
preprint cites Chen's paper only for other results and the aiXiv version
nowhere.
The preprint proves the same statement as \ploc{Lemma~6.3,
pp.~19--20}{variation.tex:130--213}, which settles the terminal and
empty-support cases of the \ploc{proof of Theorem~1.1,
p.~20}{variation.tex:215--221}, and gives the proof as its own
(\ploc{p.~19}{variation.tex:137--139}; \ploc{\S1.1, p.~3}{main.tex:183--185}),
crediting only the conclusion to Auffinger, Chen and Zeng
\sloc{AuffingerEtAl2017-SKInfiniteStepRSBZeroTemp}{Remark~7, p.~15}, who
assert it and indicate a different, perturbative route without details.
Both proofs run a boundary layer at $t=1$: the Cole--Hopf formula makes
$u(1-\epsilon,\sqrt\epsilon\,y)$ a sign average under the tilt
$e^{c\sqrt\epsilon|y+Z|}$, harmless at this scale, so $1-u^2$ is of order one
on a window of width $\sqrt\epsilon$ where $X_t$ has density of order one;
then $1-q(t)$ is at least of order $\sqrt{1-t}$, and $G<0$ near~$1$
contradicts $G\ge0$.
The case $c=0$ goes back to Chen, Handschy and Lerman
\sloc{ChenEtAl2016-EnergyLandscapeMixedEvenPSpin}{proof of Thm.~6,
pp.~14--16}, for $\gamma\equiv0$, and to El Alaoui, Montanari and Sellke
\cite{ElAlaouiEtAl2020-OptimizationMeanFieldSpinGlasses}; the extension to a
positive plateau, which behaves like $c=0$ in this layer, first appears in
the aiXiv version of Chen's paper.
\rel{same result}: as $0\in S$ (Chen's Prop.~3.2), $S$ is closed and the
minimizer unique \sloc{ChenEtAl2016-EnergyLandscapeMixedEvenPSpin}{Thm.~4,
p.~8}, every terminal plateau lies in one starting at $\max S\in S$, so
Chen's Proposition~4.5 is Lemma~6.3; the preprint replaces Chen's limits
(4.30) and (4.33), and the asymptotic (4.38), by fixed-window lower bounds
adapted from the El Alaoui--Montanari--Sellke argument, which it calls
related (\ploc{p.~19}{variation.tex:140--143}).

\paragraph{Correspondence.}
Here $\epsilon=1-t$, $h(t,y)=\E e^{c|y+\sqrt{1-t}\,Z|}$, $p_v$ is the
$N(0,v)$ density, $\rho_s$ the law of $X_s$ ($s\in(a,1)$ fixed); Chen's $u$
is the preprint's $\Phi$, and Chen's $\rho_t$ the density of $X_t$.
\begin{center}\small
\begin{xltabular}{\linewidth}{@{}>{\raggedright\arraybackslash}p{0.15\linewidth}>{\raggedright\arraybackslash}X>{\raggedright\arraybackslash}X@{}}
\toprule
 & Preprint, Lemma~6.3 & Chen \cite{Chen2026-FRSBFullIntervalSupportZeroTemp}\\
\midrule
\endfirsthead
\toprule
 & Preprint, Lemma~6.3 & Chen \cite{Chen2026-FRSBFullIntervalSupportZeroTemp}\\
\midrule
\endhead
Statement & for a minimizer, no $a<1$ and finite $c\ge0$ with $\gamma=c$ on
$(a,1)$ (\ploc{p.~19}{variation.tex:130--134}) & for the minimizer
$\gamma_\star$, no $a<1$ with $(a,1)\cap S=\varnothing$ and $a\in S$,
Prop.~4.5, p.~23; $0\in S$, Prop.~3.2, p.~8\\
\addlinespace
Gradient at scale $\sqrt\epsilon$ & for $|y|\le\sqrt\epsilon$,
$u(t,y)=\E[\operatorname{sgn}(y+\sqrt\epsilon Z)e^{c|y+\sqrt\epsilon Z|}]/
\E e^{c|y+\sqrt\epsilon Z|}$ (\ploc{p.~20}{variation.tex:187--195}) &
$\mathsf B_\varepsilon(y):=u_x(1-\varepsilon,\sqrt\varepsilon\,y)=
\E[\operatorname{sign}(y+Z)e^{m\sqrt\varepsilon|y+Z|}]/
\E e^{m\sqrt\varepsilon|y+Z|}$, (4.28)--(4.29), p.~23\\
\addlinespace
Law of $X_t$ on the plateau &
$k_{s,t}(x,y)=p_{t-s}(y-x)h(t,y)/h(s,x)$
\ploc{(6.6), p.~19}{variation.tex:157--169}; its mixture against $\rho_s$
is bounded below on $|y|\le1$ \ploc{(6.7), p.~20}{variation.tex:171--185} &
$Q=\rho_te^{-mu}$ solves the heat equation, (3.21)--(3.22), p.~10;
$\rho_{1-\varepsilon}(\sqrt\varepsilon\,y)\to\rho_1(0)>0$, (4.33), p.~24\\
\addlinespace
Layer and contradiction & $1-q(t)\ge C\sqrt{1-t}$
\ploc{(6.5), p.~19}{variation.tex:145--149}; $G<0$ near $1$, against
Lemma~6.2 (\ploc{p.~20}{variation.tex:209--212}) &
$1-\Gamma(t)\sim(4/\sqrt\pi)\,\rho_1(0)\sqrt{1-t}$, (4.38), and $G<0$ near
$1$, against (3.8), p.~25\\
\bottomrule
\end{xltabular}
\end{center}

\paragraph{How to check.}
\begin{enumerate}
\item The display for $u(t,y)$ (\ploc{p.~20}{variation.tex:187--195}) after
$y\mapsto\sqrt\epsilon\,y$ against
\sloc{Chen2026-FRSBFullIntervalSupportZeroTemp}{(4.28)--(4.29), p.~23} with
$m=c$; \ploc{(6.6)}{variation.tex:157--169} against
\sloc{Chen2026-FRSBFullIntervalSupportZeroTemp}{(3.21)--(3.22), p.~10} with
$e^{mu}=h$.
\item The end of the proof of Lemma~6.3 (\ploc{p.~20}{variation.tex:196--212})
against \sloc{Chen2026-FRSBFullIntervalSupportZeroTemp}{(4.35)--(4.38) and
Step~4, pp.~24--25}, and the whole proof against
\sloc{ChatGPT2026-FRSBFullIntervalSupportZeroTemp}{(4.19)--(4.25),
pp.~32--33}; \ploc{(6.7)}{variation.tex:171--185} and the bound
$1-u^2\ge\kappa$ against
\sloc{ElAlaouiEtAl2020-OptimizationMeanFieldSpinGlasses}{Lemma~6.12, p.~29,
and the end of the proof of Thm.~5, p.~31}, where $c=0$.
\end{enumerate}

\space%
\finding[The origin lies in the support]{The origin lies in the support: no initial interval of zero mass}{chl-origin-in-support}

\paragraph{Summary.}
Chen, Handschy and Lerman (CHL) proved in 2016 that at zero external field
the smallest point of the support of the zero-temperature Parisi measure
is~$0$ \sloc{ChenEtAl2016-EnergyLandscapeMixedEvenPSpin}{\S3.2, p.~18}, for
the mixed even $p$-spin models, SK with $\xi(s)=s^2/2$ among them. The
preprint cites their paper only for other results, such as uniqueness and
the variational conditions (\ploc{\S1.1, p.~3}{main.tex:173--175, 207--211};
\ploc{pp.~5, 18, 20}{regularity.tex:9--11; variation.tex:8--10;
order-parameter-regularity.tex:10--11}), but not for this result, which it
proves again without a citation as the initial-component case of the
\ploc{proof of Theorem~1.1, p.~20}{variation.tex:239--242}.
The argument it uses, and the insight behind it, are Auffinger and Chen's
\sloc{AuffingerChen2013-PropertiesParisiMeasures}{proof of Thm~1(i), p.~15}
for the positive-temperature analogue $0\in\operatorname{supp}\mu_P$: on an
initial interval of zero mass the slope of the self-consistency function
increases strictly, so the function, which vanishes at $0$ by symmetry,
cannot climb to the diagonal at a positive first support point~$b$, where
the contact condition bounds its slope by~$1$; the preprint cites that paper
only for its moment method (\ploc{pp.~3, 20}{main.tex:215--218;
order-parameter-regularity.tex:7--9}).
\rel{same result}: CHL's statement is the conclusion of the step, for the
same functional and class, and their overlap map $\psi_1$ is $q$ on $[0,q_h]$ (below), so
they reach the same two-fixed-point contradiction, which they close with
W.-K. Chen's fixed-point uniqueness theorem \cite[Thm~5]{NL-Chen2014} rather
than with a slope bound on the gap.
H.-B. Chen's 2026 paper \sloc{Chen2026-FRSBFullIntervalSupportZeroTemp}{Prop.~3.2,
pp.~8--9}, like its aiXiv version
\sloc{ChatGPT2026-FRSBFullIntervalSupportZeroTemp}{p.~17}, states $0\in S$
and takes it from CHL, and El Alaoui, Montanari and Sellke
\sloc{ElAlaouiEtAl2020-OptimizationMeanFieldSpinGlasses}{proof of Thm~5,
pp.~30--31} run the same contradiction over a larger, non-monotone class;
the preprint cites H.-B. Chen's paper (not its aiXiv version) and that proof
only for other content (the proof for its terminal segment,
\ploc{p.~19}{variation.tex:140--143}).

\paragraph{Correspondence.}
Here $b=\min S$ if $0\notin S$, and $q_h$ is the smallest support point of
CHL's minimizer $\gamma_h$. At $h=0$ their overlap map is
$\psi_1(s)=\E\,\partial_x\Phi_{\gamma_h}(q_h,\chi_1)\,\partial_x\Phi_{\gamma_h}(q_h,\chi_2)$,
$\E\chi_i^2=\xi'(q_h)$, $\E\chi_1\chi_2=\xi'(s)$. As
$u(t,B_t)=\E[u(b,B_b)\mid B_t]$ on $[0,b]$,
$q(s)=\E[u(b,\chi_1)u(b,\chi_2)]$ with $\Var\chi_i=b$,
$\Cov(\chi_1,\chi_2)=s$, which is $\psi_1(s)$ when $\xi'(s)=s$ and $q_h=b$.
\begin{center}\small
\begin{xltabular}{\linewidth}{@{}>{\raggedright\arraybackslash}p{0.11\linewidth}>{\raggedright\arraybackslash}X>{\raggedright\arraybackslash}X>{\raggedright\arraybackslash}X@{}}
\toprule
 & Preprint & CHL \cite{ChenEtAl2016-EnergyLandscapeMixedEvenPSpin} & Auffinger--Chen \cite{AuffingerChen2013-PropertiesParisiMeasures}\\
\midrule
\endfirsthead
\toprule
 & Preprint & CHL \cite{ChenEtAl2016-EnergyLandscapeMixedEvenPSpin} & Auffinger--Chen \cite{AuffingerChen2013-PropertiesParisiMeasures}\\
\midrule
\endhead
Statement & $0\in S$, any minimizer
(\ploc{p.~20}{variation.tex:239--242}) & $q_h=0$ at $h=0$ (p.~18; also
p.~4) & $0\in\operatorname{supp}\mu_P$, positive temperature (Thm~1(i),
p.~4)\\
\addlinespace
Gap ends & $q(0)=0$ (\ploc{Lemma~2.2, p.~7}{regularity.tex:200--208});
$q(b)=b$, $q'(b)\le1$ by \ploc{(6.4), p.~19}{variation.tex:94--97} &
$\psi_1(0)=0$ (Lemma~4, p.~11); $\psi_1(q_h)=q_h$,
$\xi''(q_h)\E(\partial_{xx}\Phi_{\gamma_h})^2\le1$ (34)--(35), p.~18 &
$\Gamma=\Gamma_{\mu_P}$ of (18), p.~8: $\Gamma(0)=0$; $\Gamma(q_m)=q_m$,
$\Gamma'(q_m)\le1$ (Thm~5, p.~9; p.~15)\\
\addlinespace
Closing step & $(q')'=\E\Phi_{xxx}^2>0$ on $(0,b)$
(\ploc{Lemma~2.3, p.~8}{regularity.tex:260, 298--307}), so $q'<q'(b)\le1$
and $q(b)<b$ & two fixed points $0$, $q_h$ of $\psi_1$, against W.-K. Chen's
Thm~5 & $\Gamma''=\gamma_{1,\mu_P}>0$ on $(0,q_m)$ (Prop.~3, pp.~8--9): one
fixed point on $[0,q_m]$\\
\bottomrule
\end{xltabular}
\end{center}

\paragraph{How to check.}
\begin{enumerate}
\item The initial-component case (\ploc{p.~20}{variation.tex:239--242}) and
the $c=0$ clause of \ploc{Lemma~2.3}{regularity.tex:260} (proof on
\ploc{p.~8}{regularity.tex:298--307}) against
\sloc{AuffingerChen2013-PropertiesParisiMeasures}{proof of Thm~1(i), p.~15}
and \sloc{AuffingerChen2013-PropertiesParisiMeasures}{Prop.~3, (20)--(22),
pp.~8--9} with $\mu([0,u])=0$, $\xi'''=0$.
\item The same step against
\sloc{ChenEtAl2016-EnergyLandscapeMixedEvenPSpin}{\S3.2, p.~18}: (34)--(35)
are $q(b)=b$, $q'(b)\le1$ of \ploc{(6.4)}{variation.tex:94--97}; and
\sloc{ChenEtAl2016-EnergyLandscapeMixedEvenPSpin}{(5)--(6), Thm~4, p.~8} at
$h=0$ against \ploc{(1.2)--(1.3)}{main.tex:75--85}.
\item \sloc{Chen2026-FRSBFullIntervalSupportZeroTemp}{end of the proof of
Prop.~3.2, p.~9} and \sloc{ChatGPT2026-FRSBFullIntervalSupportZeroTemp}{p.~17}:
$0\in S$, taken from CHL.
\end{enumerate}

\space%
\finding[Full support verifies the no-overlap-gap assumption]{Full support verifies the no-overlap-gap assumption of the optimization algorithm}{chen-eams-hypothesis-verified}

\paragraph{Summary.}
Chen's paper \sloc{Chen2026-FRSBFullIntervalSupportZeroTemp}{\S1, p.~4}
already observes that its zero-temperature full-support theorem verifies, for
SK, the no-overlap-gap assumption of El Alaoui, Montanari and Sellke, under
which their message-passing algorithm comes within any $\epsilon>0$ of the
optimum. The preprint cites Chen's paper only for the full-support theorem and
proof steps: its \ploc{Corollary~7.3, p.~25}{algorithm-application.tex:13--34}
and the announcements of it on
\ploc{pp.~2, 3 and~25}{main.tex:145--147, 225--230; algorithm-application.tex:4--5}
name or cite only El Alaoui, Montanari and Sellke, for the assumption and the
corollary. The shared approach is one deduction: full support makes the
minimizer strictly increasing on $[0,1)$, which is exactly their Assumption~2
for SK, so their Corollary~2.2 applies; they posed the assumption for this use
and expected it to hold for SK
\sloc{ElAlaouiEtAl2020-OptimizationMeanFieldSpinGlasses}{p.~8}, and the
observation that the zero-temperature full-support theorem settles it for SK is
already in Chen's paper (the theorem itself is
finding~\ref{f:chen-same-main-theorem}). The relationship is
\rel{same result}, since both papers draw the same conclusion by this
deduction; to apply Corollary~2.2, the preprint adds an independent diagonal
that puts $H_n$ in its Gaussian input format, carrying over the guarantee and
the cost, and takes the operation count $C(\epsilon)n^2$ from Theorem~3 of El
Alaoui, Montanari and Sellke. Corollary~7.3 feeds no other result of the
preprint, hence \emph{peripheral}.

\paragraph{Correspondence.}
\begin{center}\small
\begin{xltabular}{\linewidth}{@{}>{\raggedright\arraybackslash}p{0.17\linewidth}>{\raggedright\arraybackslash}X>{\raggedright\arraybackslash}X@{}}
\toprule
 & Preprint & Chen \cite{Chen2026-FRSBFullIntervalSupportZeroTemp}\\
\midrule
\endfirsthead
\toprule
 & Preprint & Chen \cite{Chen2026-FRSBFullIntervalSupportZeroTemp}\\
\midrule
\endhead
Strict increase & $\operatorname{supp}\mu_\gamma=[0,1)$, so $\gamma(b)>\gamma(a)$ for $0\le a<b<1$ (\ploc{Theorem~1.1, p.~2}{main.tex:119--124}) & $\operatorname{supp}_{[0,1)}\nu_\star=[0,1)$ (Theorem~1.3, p.~3), so $\gamma_\star(t)-\gamma_\star(s)=\nu_\star((s,t])>0$ for $0\le s<t<1$ (p.~4)\\
\addlinespace
Assumption & Assumption~2 of El Alaoui et al.\ for $\xi(t)=t^2/2$ (\ploc{proof of Corollary~7.3}{algorithm-application.tex:31--33}) & Assumption~2 of El Alaoui et al.\ for SK (p.~4)\\
\addlinespace
Conclusion & their Corollary~2.2: $\Pr\{H_n(\sigma^{\rm alg})\ge\max_\sigma H_n(\sigma)-\epsilon n\}\to1$ (\ploc{Corollary~7.3}{algorithm-application.tex:15--26}) & their Corollary~2.2: energy within any fixed $\varepsilon>0$ of the limiting optimum (p.~4)\\
\addlinespace
\bottomrule
\end{xltabular}
\end{center}

\paragraph{How to check.}
\begin{enumerate}
\item The proof of \ploc{Corollary~7.3, p.~25}{algorithm-application.tex:31--34}
against \sloc{Chen2026-FRSBFullIntervalSupportZeroTemp}{p.~4}, the paragraph
beginning ``Full support is also relevant to optimization algorithms,'' and
both against
\sloc{ElAlaouiEtAl2020-OptimizationMeanFieldSpinGlasses}{Assumption~2,
Corollary~2.2 and (2.8), p.~8}.
\item The paragraphs on \ploc{pp.~2--3}{main.tex:133--230} and
\ploc{\S7.1, pp.~25--26}{algorithm-application.tex:1--59}: Chen's paper is
cited for Theorem~1.3 and for proof steps, but not for the algorithmic
application.
\end{enumerate}

\space%

\begin{thebibliography}{99}
\small
\bibitem{Chen2026-FRSBFullIntervalSupportZeroTemp} Hong-Bin Chen, FRSB in the SK spin glass: convergence to full-interval support at zero temperature, arXiv preprint, 2026. arXiv:2607.18032v1.
\bibitem{ChatGPT2026-FRSBFullIntervalSupportZeroTemp} ChatGPT, FRSB in the SK Spin Glass: Convergence to Full-Interval Support at Zero Temperature, aiXiv preprint, 2026. aiXiv:aixiv.260716.000001.
\bibitem{Zhou2025-FRSBSKLowTemperature} Yuxin Zhou, Existence of full replica symmetry breaking for the Sherrington-Kirkpatrick model at low temperature, Comm. Pure Appl. Math. 79(7) (2026) 1746-1770. arXiv:2504.00269v2.
\bibitem{Lopatto2026-FullRSBSK} Patrick Lopatto, Full replica symmetry breaking in the Sherrington-Kirkpatrick model, arXiv preprint, 2026. arXiv:2607.11756v3.
\bibitem{Zhou2026-GardnerTransitionIsingPurePSpinII} Yuxin Zhou, On the Gardner transition in the Ising pure p-spin glass II, arXiv preprint, 2026. arXiv:2608.06523v1.
\bibitem{AuffingerEtAl2017-SKInfiniteStepRSBZeroTemp} Antonio Auffinger, Wei-Kuo Chen, Qiang Zeng, The SK model is infinite step replica symmetry breaking at zero temperature, Comm. Pure Appl. Math. (2020). arXiv:1703.06872v2.
\bibitem{ElAlaouiEtAl2020-OptimizationMeanFieldSpinGlasses} Ahmed El Alaoui, Andrea Montanari, Mark Sellke, Optimization of mean-field spin glasses, Ann. Probab. 49 (2021). arXiv:2001.00904v1.
\bibitem{ChenEtAl2016-EnergyLandscapeMixedEvenPSpin} Wei-Kuo Chen, Madeline Handschy, Gilad Lerman, On the energy landscape of the mixed even p-spin model, Probab. Theory Related Fields (2018). arXiv:1609.04368v4.
\bibitem{AuffingerChen2013-PropertiesParisiMeasures} Antonio Auffinger, Wei-Kuo Chen, On properties of Parisi measures, Probab. Theory Related Fields (2015). arXiv:1303.3573v1.
\bibitem{OAI-Preprint} OpenAI, \emph{Full support of the zero-temperature Sherrington--Kirkpatrick order parameter}, OpenAI Math Release preprint, September 27, 2026, \url{https://github.com/openai/math}.
\bibitem{NL-Chen2014} Wei-Kuo Chen, Chaos in the mixed even-spin models, Comm. Math. Phys. 328 (2014), 867--901, doi:10.1007/s00220-014-1888-1
\end{thebibliography}
\space%

\end{document}
